\section{Virtual Tree}
\whenuse{Cada consulta marca pocos vértices de un árbol grande y sólo importa
la unión de los caminos que los conectan.}

El árbol virtual comprime el árbol original conservando únicamente los
vértices marcados y los LCA necesarios para conectarlos. Una arista virtual
$(u,v)$ representa el camino completo desde el ancestro $u$ hasta $v$ en el
árbol original.

Por ejemplo, si los marcados son $a,b,c$, puede ser necesario agregar
$LCA(a,b)$ para conectarlos. No se agregan los vértices intermedios que no
cambian la estructura de los caminos.

Requiere precomputar \texttt{tin}, \texttt{tout} y LCA. La función devuelve
los vértices virtuales y sus aristas \texttt{(padre,hijo)}.

\complexity{Por consulta $O(k\log k)$ tiempo y $O(k)$ memoria; para $k$
marcados se crean como máximo $2k-1$ vértices.}

\begin{lstlisting}[style=cpp]
template<class LCAFunction>
pair<vector<int>,vector<pair<int,int>>> build_virtual_tree(
        vector<int> marked,const vector<int>& tin,const vector<int>& tout,
        LCAFunction get_lca) {
    auto ancestor=[&](int a,int b) {
        return tin[a]<=tin[b]&&tout[b]<=tout[a];
    };
    auto by_tin=[&](int a,int b) { return tin[a]<tin[b]; };
    sort(marked.begin(),marked.end(),by_tin);
    marked.erase(unique(marked.begin(),marked.end()),marked.end());
    int original_size=marked.size();
    for(int i=1;i<original_size;++i)
        marked.push_back(get_lca(marked[i-1],marked[i]));
    sort(marked.begin(),marked.end(),by_tin);
    marked.erase(unique(marked.begin(),marked.end()),marked.end());
    vector<pair<int,int>> edges;
    vector<int> st;
    for(int vertex:marked) {
        while(!st.empty()&&!ancestor(st.back(),vertex)) st.pop_back();
        if(!st.empty()) edges.push_back({st.back(),vertex});
        st.push_back(vertex);
    }
    return {marked,edges};
}

// Uso
auto [vertices,edges]=build_virtual_tree(
    marked,tin,tout,[&](int a,int b){ return lca.query(a,b); });
for(auto [parent,child]:edges)
    cout<<parent<<" "<<child<<" "
        <<depth[child]-depth[parent]<<'\n';
\end{lstlisting}

Si el árbol original tiene pesos, reemplazar la diferencia de profundidades por:

\[
    distRoot[child]-distRoot[parent].
\]

\warning{Los LCA agregados no necesariamente estaban marcados originalmente.
Guardar aparte los marcados si la DP necesita distinguirlos. Incluir la raíz
original explícitamente sólo si la consulta necesita que forme parte del árbol
virtual.}